% generated by roam.py handbook export -- не редактировать, править заметки
% 20261014115449-глобальная_сходимость_метода_ньютона.org (body 51dd588e5f48)

\begin{theorem}[{Глобальная сходимость метода Ньютона}]\label[theorem]{note:CC41F7D2-EECB-46E4-817A-756E7CB9B7C0}
  Пусть

  \begin{itemize}
    \item $f \in C_L^{1,1}$ (\hyperref[note:7DD1C7AB-1A49-4675-9F21-5A16C717F670]{Функции с липшицевым градиентом}),
    \item $f \in C^2$ (нужен \hyperref[note:592AD816-1D80-48E6-8399-8524FBC1E8F3]{гессиан})
    \item $f$~--- $\mu$-\hyperref[note:22E546FC-186F-4A67-85AA-2C85DCEF9021]{сильно выпуклая}
    \item $x_{k+1} = x_k - \alpha_k(\nabla^2 f(x_k))^{-1}\nabla f(x_k)$ (направление по \hyperref[note:C0ED5501-CFC6-42BA-A53A-1B346CBA545D]{Ньютону})
    \item $\alpha_k$~--- выполняет \hyperref[note:5330B23D-95DD-4D5D-B4A3-63EBADB1DB7D]{Условия Армихо-Вульфа} с $0 < c_1 < c_2 < 1$
  \end{itemize}

  Тогда $\nabla f(x_k)\rightarrow0$, при этом $\min_{1\le i \le k} \|\nabla f(x_i)\| \le \frac{C}{\sqrt{k}}$
\end{theorem}

\begin{proof}
  Пусть $\cos \theta_k$~--- \hyperref[note:E4AE9B25-BDB9-4A54-8CCB-09809F20B79D]{косинус угла между производной и направлением спуска}.
  Здесь $d_k = -(\nabla^2 f(x_k))^{-1} \nabla f(x_k)$. Так как $\nabla^2 f(x_k) \ge \mu I > 0$, гессиан обратим и $(\nabla^2 f(x_k))^{-1} > 0$, поэтому при $\nabla f(x_k) \ne 0$
  \[
  \left\langle \nabla f(x_k); d_k \right\rangle = - \left\langle \nabla f(x_k); (\nabla^2 f(x_k))^{-1} \nabla f(x_k) \right\rangle < 0
  \]
  -- это направление спуска. В $\cos^2 \theta_k$ знак не важен, поэтому дальше минус опускаем.
  Нам достаточно доказать, что
  \[
  \cos^2 \theta_k = \left( \frac{\left\langle \nabla f(x_k); d_k \right\rangle}{\|\nabla f(x_k)\| \cdot \| d_k \|} \right)^2 > c
  \]
  Это следует из \hyperref[note:93FC9CC4-79FA-4CEA-9B10-9DBFB2C07850]{оценки на шаг Армихо Вульфа} и \hyperref[note:36F6E1BA-41DB-460D-A9C6-B980621727D4]{сходимости градиентного спуска с шагом Армихо-Вульфа}

  Давайте заметим, что
  \[
  A = A^T \Rightarrow \lambda_{\min}(A) \|x\|^2 \le \left\langle Ax;x \right\rangle \le \lambda_{\max}(A) \|x\|^2
  \]
  Это достаточно очевидно, см. \hyperref[note:3BE0BBCE-AAB0-4CB4-8FA6-45277C0889CF]{спектральную теорему}

  Тогда
  \[
  \cos^2 \theta_k = \left( \frac{\left\langle \nabla f(x_k); (\nabla^2 f(x_k))^{-1} \nabla f(x_k) \right\rangle}{\|\nabla f(x_k) \| \cdot \| (\nabla^2 f(x_k))^{-1} \nabla f(x_{k})\|} \right)^2
  \]
  Оценим числитель снизу, заметим что он подходит под наше замечание про оценку $\left\langle Ax;x \right\rangle$,
  так как Гессиан~--- симметричен;
  \[
  \left\langle \nabla f(x_k); (\nabla^2 f(x_k))^{-1} \nabla f(x_k) \right\rangle \ge
  \lambda_{\min}\left ((\nabla^2 f(x_k))^{-1}\right) \|\nabla f(x_k)\|^2
  \]

  Оценим произведение нормы матрицы на вектор
  \[
  \| (\nabla^2 f(x_k))^{-1} \nabla f(x_{k})\| \le \|(\nabla^2 f(x_k))^{-1}\| \cdot \|\nabla f(x_k)\|
  \le \lambda_{\max}\left ((\nabla^2 f(x_k))^{-1} \right) \|\nabla f(x_k)\|
  \]
  Последнее неравенство: $(\nabla^2 f(x_k))^{-1} > 0$, поэтому ее спектральная норма равна $\lambda_{\max}$

  Тогда вся дробь
  \begin{multline*}
  \frac{\left\langle \nabla f(x_k); (\nabla^2 f(x_k))^{-1} \nabla f(x_k) \right\rangle}{\|\nabla f(x_k) \| \cdot \| (\nabla^2 f(x_k))^{-1} \nabla f(x_{k})\|} \\
  \ge \frac{\lambda_{\min}\left ((\nabla^2 f(x_k))^{-1}\right) \|\nabla f(x_k)\|^2
  }{
  \|\nabla f(x_k)\| \lambda_{\max}\left ((\nabla^2 f(x_k))^{-1} \right) \|\nabla f(x_k)\|
  } = \frac{\lambda_{\min}\left ((\nabla^2 f(x_k))^{-1}\right)}{\lambda_{\max}\left ((\nabla^2 f(x_k))^{-1} \right)}
  \end{multline*}

  Давайте воспользуемся \hyperref[note:94DADA28-87E4-4E36-B289-8415BCE750FC]{критерием сильной выпуклости в терминах вторых производных} и
  \hyperref[note:DAA0DAE4-4B48-4A2F-9459-54790B6B129F]{оценкой нормы гессиана сверху}, тогда мы получим такие неравенства
  \[
  \lambda_{\min}(\nabla^2 f(x)) \ge \mu \quad \lambda_{\max}(\nabla^2 f(x)) \le L
  \]
  Тогда
  \[
  \lambda_{\min}\left ((\nabla^2 f(x_k))^{-1}\right) \ge \frac{1}{L}\quad
  \lambda_{\max}\left ((\nabla^2 f(x_k))^{-1} \right) \le \frac{1}{\mu}
  \]
  Итого
  \[
  \cos^2 \theta_k \ge \left( \frac{\mu}{L} \right)^2
  \]

  Тогда по \hyperref[note:36F6E1BA-41DB-460D-A9C6-B980621727D4]{сходимости градиентного спуска с шагом Армихо-Вульфа} с $C = \frac{\mu}{L}$ (у сильно выпуклой функции $f^{*}$ существует):
  \[
  \min_{0\le k\le n} \|\nabla f(x_k)\| \le \frac{L}{\mu} \sqrt{\frac{L (f(x_0) - f^{*})}{(n+1)c_1(1-c_2)}}
  \]

  \begin{remark}
    Сходятся и сами точки: для $\mu$-сильно выпуклой $f$ верно $\|\nabla f(x)\| \ge \mu \|x - x^{*}\|$, поэтому из $\|\nabla f(x_k)\| \rightarrow 0$ следует $x_k \rightarrow x^{*}$.
  \end{remark}
\end{proof}
